% !TeX root = ../main.tex
\begin{frame}{器件验证与零泊松比支架}
  \begin{SlideCanvas}
    \LiteratureTimeAxis
    \LiteratureTimeNode{190}{2021}
    \LiteratureTimeNode{480}{2026}
    \LiteratureTimeNode{770}{2026}
    \PointAt{55}{153}{270}{Wang等：PCL}
    \PointAt{345}{153}{270}{Gupta等：金属原型}
    \PointAt{635}{153}{270}{Yasmin等：镍钛}
    \CenteredPicAt{190}{275}{260}{180}{assets/wang2021_length.pdf}
    \CenteredPicAt{480}{275}{260}{180}{assets/gupta2026_fig10.jpg}
    \CenteredPicAt{770}{275}{260}{180}{assets/zero2026_geometry.png}
    \DetailAt{55}{383}{270}{打印PCL零泊松比支架，通过径向扩张试验观察轴向长度稳定性，并评价径向力学性能。}
    \DetailAt{345}{383}{270}{飞秒激光切割与抛光制备金属原型，径向强度提高，同时出现轴向伸长。}
    \DetailAt{635}{383}{270}{NiTi压握／释放建模与打印原型；15--25°构型在50\%压握终态报告$|\mathrm{FS}|\leq0.12\%$。}
    \Sources{\RefWang，图11(b)；\RefGupta，图10；\RefYasmin，图4。\\年份按卷期，节点间距示意；各项结果按其材料与加载条件解释。}
  \end{SlideCanvas}
  \BICConclusionBand{零泊松比与径向力已有共同研究，响应轨迹定制构成\BICBandEmph{进一步的设计问题}}
  \note{Wang图11(b)为原始、最大扩张及回弹状态的长度统计，保留完整柱图、纵轴与图例，不把其纵轴局部区间改绘成零起点。Gupta图10为机械测试装置；摘要中径向强度3.22对1.26N/mm、短缩率-9.44\%对+2.49\%属于同研究比较，负短缩表示伸长。Yasmin图4展示箭头单胞、阵列及圆管；长度指标针对t=0.50mm、l=4.50mm、15--25度构型的50\%压握终态，不是释放全过程最大偏差。NiTi成形件相变终温高于体温以及端部约束等研究范围在挑战页讨论。不同材料、尺寸和分支不做统一绝对值排名。}
\end{frame}
